\chapter{2006年江西卷（理）}
\section{单选题}
\begin{problem}
已知集合 \(M = \left\{x \mid \dfrac{x}{(x - 1)^3} \geqslant 0\right\}\)，\(N = \{y \mid y = 3x^2 + 1, x \in \R\}\)，则 \(M \cap N\) 等于（\quad）

\choices
    {\(\varnothing\)}
    {\(\{x\mid x\geqslant 1\}\)}
    {\(\{x\mid x>1\}\)}
    {\(\{x\mid x\geqslant 1 \text{\ 或\ } x<0\}\)}
\end{problem}

\begin{answer}
C
\end{answer}

\begin{solution}
由分子、分母的零点及定义域得，\(M=(-\infty,0]\cup(1,+\infty)\). 又因为 \(y=3x^2+1\geqslant1\)，所以 \(N=[1,+\infty)\). 因此 \(M\cap N=(1,+\infty)\)，选 C.
\end{solution}

\begin{problem}
已知复数 \(z\) 满足 \((\sqrt{3} + 3\i)z = 3\i\)，则 \(z\) 等于（\quad）

\choices
    {\(\frac{3}{2} - \frac{\sqrt{3}}{2}\i\)}
    {\(\frac{3}{4} - \frac{\sqrt{3}}{4}\i\)}
    {\(\frac{3}{2} + \frac{\sqrt{3}}{2}\i\)}
    {\(\frac{3}{4} + \frac{\sqrt{3}}{4}\i\)}
\end{problem}

\begin{answer}
D
\end{answer}

\begin{solution}
将等式两边同除以 \(\sqrt{3}+3\i\)，并乘以共轭复数，得
\(z=\dfrac{3\i}{\sqrt{3}+3\i}=\dfrac{3\i(\sqrt{3}-3\i)}{(\sqrt{3})^2+3^2}=\dfrac{9+3\sqrt{3}\i}{12}=\dfrac{3}{4}+\dfrac{\sqrt{3}}{4}\i\). 故选 D.
\end{solution}

\begin{problem}
若 \(a > 0\)，\(b > 0\)，则不等式 \(-b < \frac{1}{x} < a\) 等价于（\quad）

\choices
    {\(-\frac{1}{b} < x < 0 \text{\ 或\ } 0 < x < \frac{1}{a}\)}
    {\(-\frac{1}{a} < x < \frac{1}{b}\)}
    {\(x < -\frac{1}{a}\) 或 \(x > \frac{1}{b}\)}
    {\(x < -\frac{1}{b}\) 或 \(x > \frac{1}{a}\)}
\end{problem}

\begin{answer}
D
\end{answer}

\begin{solution}
由于 \(x\neq0\)，分两种情况讨论. 当 \(x>0\) 时，\(1/x>0>-b\)，且 \(1/x<a\) 等价于 \(x>1/a\). 当 \(x<0\) 时，\(1/x<a\) 自动成立，而 \(-b<1/x\) 乘以负数 \(x\) 后等价于 \(x<-1/b\). 所以原不等式等价于 \(x<-1/b\) 或 \(x>1/a\)，选 D.
\end{solution}

\begin{problem}
设 \(O\) 为坐标原点，\(F\) 为抛物线 \(y^{2} = 4x\) 的焦点，\(A\) 为抛物线上的一点，若 \( \overrightarrow{OA} \cdot \overrightarrow{AF} = -4 \)，则点 \(A\) 的坐标为（\quad）

\choices
    {\((2,\pm 2\sqrt{2})\)}
    {\((1,\pm 2)\)}
    {\((1,2)\)}
    {\((2,2\sqrt{2})\)}
\end{problem}

\begin{answer}
B
\end{answer}

\begin{solution}
抛物线 \(y^2=4x\) 的焦点为 \(F(1,0)\). 设 \(A(x,y)\)，则 \(y^2=4x\)，且
\(\overrightarrow{OA}\cdot\overrightarrow{AF}=(x,y)\cdot(1-x,-y)=x-x^2-y^2=-x^2-3x\).
由题设 \(-x^2-3x=-4\)，即 \((x-1)(x+4)=0\). 抛物线上 \(x\geqslant0\)，故 \(x=1\)，从而 \(y=\pm2\). 所以 \(A=(1,\pm2)\)，选 B.
\end{solution}

\begin{problem}
对于 \(\R\) 上可导的任意函数 \(f(x)\)，若满足 \((x - 1)f'(x) \geqslant 0\)，则必有（\quad）

\choices
    {\(f(0) + f(2) < 2f(1)\)}
    {\(f(0) + f(2)\leqslant 2f(1)\)}
    {\(f(0) + f(2)\geqslant 2f(1)\)}
    {\(f(0) + f(2) > 2f(1)\)}
\end{problem}

\begin{answer}
C
\end{answer}

\begin{solution}
当 \(x<1\) 时，题设给出 \(f'(x)\leqslant0\)，所以 \(f\) 在 \((-\infty,1]\) 上单调不增；当 \(x>1\) 时，\(f'(x)\geqslant0\)，所以 \(f\) 在 \([1,+\infty)\) 上单调不减. 于是 \(f(0)\geqslant f(1)\)，\(f(2)\geqslant f(1)\)，从而 \(f(0)+f(2)\geqslant2f(1)\)，选 C.
\end{solution}

\begin{problem}
若不等式 \(x^{2} + ax + 1 \geqslant 0\) 对一切 \(x \in \left(0, \frac{1}{2}\right]\) 成立，则 \(a\) 的最小值为（\quad）

\choices
    {0}
    {\(-2\)}
    {\(-\frac{5}{2}\)}
    {\(-3\)}
\end{problem}

\begin{answer}
C
\end{answer}

\begin{solution}
对任意 \(x\in(0,1/2]\)，不等式等价于 \(a\geqslant-x-1/x\). 函数 \(-x-1/x\) 在 \((0,1/2]\) 上的导数为 \(-1+1/x^2>0\)，故其最大值在 \(x=1/2\) 处取得，等于 \(-1/2-2=-5/2\). 因此 \(a\) 的最小值为 \(-5/2\)，选 C.
\end{solution}

\begin{problem}
已知等差数列 \(\{a_{n}\}\) 的前 \(n\) 项和为 \(S_{n}\). 若 \(\overrightarrow{OB} = a_{1}\overrightarrow{OA} + a_{200}\overrightarrow{OC}\)，且 \(A\)，\(B\)，\(C\) 三点共线（该直线不过点 \(O\)），则 \(S_{200}\) 等于（\quad）

\choices
    {100}
    {101}
    {200}
    {201}
\end{problem}

\begin{answer}
A
\end{answer}

\begin{solution}
因为 \(A\)，\(B\)，\(C\) 共线且该直线不过原点，所以 \(\overrightarrow{OA}\) 与 \(\overrightarrow{OC}\) 不共线. 点 \(B\) 在直线 \(AC\) 上，故其位置向量表示中两个系数之和必须为 \(1\)，由题设得到 \(a_1+a_{200}=1\). 等差数列的前 \(200\) 项和为首项与末项平均值乘以项数，即 \(S_{200}=100(a_1+a_{200})=100\)，选 A.
\end{solution}

\begin{problem}
在 \((x - \sqrt{2})^{2006}\) 的二项展开式中，含 \(x\) 的奇次幂的项之和为 \(S\)，当 \(x = \sqrt{2}\) 时，\(S\) 等于（\quad）

\choices
    {\(2^{3008}\)}
    {\(-2^{3008}\)}
    {\(2^{3009}\)}
    {\(-2^{3009}\)}
\end{problem}

\begin{answer}
B
\end{answer}

\begin{solution}
设 \(P(x)=(x-\sqrt{2})^{2006}\)，将展开式按 \(x\) 的幂次分为奇次幂部分 \(S(x)\) 与偶次幂部分. 则 \(S(x)=\dfrac{P(x)-P(-x)}{2}\). 代入 \(x=\sqrt{2}\)，有 \(P(\sqrt{2})=0\)，且 \(P(-\sqrt{2})=(-2\sqrt{2})^{2006}=2^{3009}\). 因此 \(S=\dfrac{0-2^{3009}}{2}=-2^{3008}\)，选 B.
\end{solution}

\begin{problem}
\(P\) 为双曲线 \(\frac{x^2}{9} - \frac{y^2}{16} = 1\) 的右支上一点，\(M\)，\(N\) 分别是圆 \((x + 5)^2 + y^2 = 4\) 和 \((x - 5)^2 + y^2 = 1\) 上的点，则 \(|PM| - |PN|\) 的最大值为（\quad）

\choices
    {6}
    {7}
    {8}
    {9}
\end{problem}

\begin{answer}
D
\end{answer}

\begin{solution}
双曲线的焦点为 \((-5,0)\)、\((5,0)\)，且右支上任意点 \(P\) 满足 \(|P(-5,0)|-|P(5,0)|=2a=6\). 由三角不等式，\(M\) 在半径为 \(2\) 的圆上时 \(|PM|\leqslant|P(-5,0)|+2\)，\(N\) 在半径为 \(1\) 的圆上时 \(|PN|\geqslant|P(5,0)|-1\). 因此
\(|PM|-|PN|\leqslant|P(-5,0)|-|P(5,0)|+3=9\). 由双曲线方程且 \(P\) 在右支上，设 \(P=(x,y)\)，则 \(x\geqslant3\)，并且
\(|P(5,0)|^2=(x-5)^2+y^2=\dfrac{(5x-9)^2}{9}\)，故 \(|P(5,0)|\geqslant2>1\). 在从 \(P\) 经过左焦点的射线上取距左焦点为 \(2\) 的点 \(M\)，在右焦点与 \(P\) 的线段上取距右焦点为 \(1\) 的点 \(N\)，即可同时取到上述两个等号，故最大值为 \(9\)，选 D.
\end{solution}

\begin{problem}
将 7 个人（含甲，乙）分成三个组，一组 3 人，另两组各 2 人，不同的分组数为 \(a\)，甲，乙分在同一组的概率为 \(p\)，则 \(a\)，\(p\) 的值分别为（\quad）

\choices
    {\(a = 105, p = \frac{5}{21}\)}
    {\(a = 105\)，\(p = \frac{4}{21}\)}
    {\(a = 210\)，\(p = \frac{5}{21}\)}
    {\(a = 210\)，\(p = \frac{4}{21}\)}
\end{problem}

\begin{answer}
A
\end{answer}

\begin{solution}
先计数. 先选出三人组，再将剩下四人分成两个无序二人组，故 \(a=\mathrm C_7^3\cdot3=105\).

甲、乙同在三人组时，第三人有 \(5\) 种选法，剩余四人有 \(3\) 种配对法，共 \(15\) 种；甲、乙同在二人组时，三人组从其余五人中选出，共 \(\mathrm C_5^3=10\) 种. 因此有利分组数为 \(25\)，\(p=25/105=5/21\). 故选 A.
\end{solution}

\begin{problem}
如图，在四面体 \(ABCD\) 中，截面 \(AEF\) 经过四面体的内切球（与四个面都相切的球）球心 \(O\)，且与 \(BC\)，\(DC\) 分别截于 \(E\)，\(F\). 如果截面将四面体分为体积相等的两部分，设四棱锥 \(A-BEFD\) 与三棱锥 \(A-EFC\) 的表面积分别为 \(S_1\) 和 \(S_2\)，则必有（\quad）

\bitmapfigure[max width=0.42\linewidth]{img/2006/jiangxi_science/q11_fig1.png}

\choices
    {\(S_{1} < S_{2}\)}
    {\(S_{1} > S_{2}\)}
    {\(S_{1} = S_{2}\)}
    {\(S_{1}\) 和 \(S_{2}\) 的大小不能确定}
\end{problem}

\begin{answer}
C
\end{answer}

\begin{solution}
记截面三角形 \(AEF\) 的面积为 \(T\)，内切球半径为 \(r\). 以球心 \(O\) 为顶点，把四棱锥 \(A-BEFD\) 分解为它与各个原四面体面片组成的棱锥. 每个原面片所在平面到 \(O\) 的距离均为 \(r\)，而公共截面 \(AEF\) 所在平面经过 \(O\)，距离为 \(0\)，所以
\(V_1=\dfrac r3(S_1-T)\). 同理，\(V_2=\dfrac r3(S_2-T)\). 题设 \(V_1=V_2\)，故 \(S_1-T=S_2-T\)，即 \(S_1=S_2\)，选 C.
\end{solution}

\begin{problem}
某地一年内的气温 \(Q(t)\)（单位：\(^{\circ} \mathrm{C}\)）与时间 \(t\)（月份）之间的关系如图所示，已知该年的平均气温为 \(10^{\circ} \mathrm{C}\). 令 \(C(t)\) 表示时间段 \([0, t]\) 的平均气温，\(C(t)\) 与 \(t\) 之间的函数关系用下列图象表示，则正确的是（\quad）

\bitmapfigure[max width=0.42\linewidth]{img/2006/jiangxi_science/q12_qgraph.png}

\choices
    {\choicebitmap[width=0.15\paperwidth]{img/2006/jiangxi_science/q12_optA.png}}
    {\choicebitmap[width=0.15\paperwidth]{img/2006/jiangxi_science/q12_optB.png}}
    {\choicebitmap[width=0.15\paperwidth]{img/2006/jiangxi_science/q12_optC.png}}
    {\choicebitmap[width=0.15\paperwidth]{img/2006/jiangxi_science/q12_optD.png}}
\end{problem}

\begin{answer}
A
\end{answer}

\begin{solution}
对 \(t>0\)，有 \(C(t)=\dfrac1t\int_0^tQ(u)\,\mathrm du\)，从而
\(C'(t)=\dfrac{Q(t)-C(t)}{t}\). 因此 \(Q(t)>C(t)\) 时 \(C(t)\) 上升，\(Q(t)<C(t)\) 时 \(C(t)\) 下降. 由题图可见，初段 \(Q(t)<C(t)\)，中段 \(Q(t)>C(t)\)，末段又有 \(Q(t)<C(t)\)，所以 \(C(t)\) 先下降、后上升、再下降；又已知 \(C(12)=10\). 综上，故选 A.
\end{solution}

\section{填空题}
\begin{problem}
数列 \( \left\{\frac{1}{4n^{2}-1}\right\} \) 的前 \(n\) 项和为 \( S_{n} \)，则 \(\displaystyle\lim_{n\to\infty} S_{n} = \)\fillinblank{}.
\end{problem}

\begin{answer}
\(\frac{1}{2}\)
\end{answer}

\begin{solution}
将通项作部分分式分解，得
\(\dfrac{1}{4k^2-1}=\dfrac12\left(\dfrac{1}{2k-1}-\dfrac{1}{2k+1}\right)\). 所以
\(S_n=\dfrac12\left(1-\dfrac{1}{2n+1}\right)=\dfrac{n}{2n+1}\)，从而 \(\displaystyle\lim_{n\to\infty}S_n=\dfrac12\).
\end{solution}

\begin{problem}
设 \( f(x)=\log_{3}(x+6) \) 的反函数为 \( f^{-1}(x) \)，若 \( [f^{-1}(m)+6]\cdot[f^{-1}(n)+6]=27 \)，则 \( f(m+n)= \)\fillinblank{}.
\end{problem}

\begin{answer}
2
\end{answer}

\begin{solution}
由 \(y=\log_3(x+6)\) 得 \(x=3^y-6\)，所以 \(f^{-1}(x)=3^x-6\). 题设条件化为 \(3^m\cdot3^n=27=3^3\)，故 \(m+n=3\). 于是 \(f(m+n)=\log_3(3+6)=2\).
\end{solution}

\begin{problem}
如图，在直三棱柱 \(ABC - A_{1}B_{1}C_{1}\) 中，底面为直角三角形，\(\angle ACB = 90^{\circ}\)，\(AC = 6\)，\(BC = CC_{1} = \sqrt{2}\)，\(P\) 是 \(BC_{1}\) 上一动点，则 \(CP + PA_{1}\) 的最小值为\fillinblank{}.

\bitmapfigure[max width=0.42\linewidth]{img/2006/jiangxi_science/q15_fig1.png}
\end{problem}

\begin{answer}
\(5\sqrt{2}\)
\end{answer}

\begin{solution}
沿公共边 \(BC_1\) 将三角形 \(CBC_1\) 展开到平面 \(A_1BC_1\) 的另一侧，设展开后的点 \(C\) 为 \(C'\). 对任意 \(P\in BC_1\)，展开保持长度，故 \(CP+PA_1=C'P+PA_1\geqslant C'A_1\)，等号当且仅当 \(P\) 为直线段 \(C'A_1\) 与 \(BC_1\) 的交点.

因为 \(BC=CC_1=\sqrt{2}\) 且 \(\angle BCC_1=90^\circ\)，所以 \(BC_1=2\). 在平面 \(A_1BC_1\) 中，\(A_1C_1=AC=6\)，\(A_1B^2=AB^2+AA_1^2=(6^2+2)+2=40\). 取平面坐标 \(B=(0,0)\)，\(C_1=(2,0)\)，则由 \(A_1B=2\sqrt{10}\)、\(A_1C_1=6\) 得 \(A_1=(2,6)\)；又展开后的 \(C'=(1,-1)\). 因此 \(C'A_1=\sqrt{1^2+7^2}=5\sqrt{2}\). 线段 \(C'A_1\) 与 \(BC_1\) 的交点坐标为 \((8/7,0)\)，确实在线段 \(BC_1\) 上，故等号可以取得，最小值为 \(5\sqrt{2}\).
\end{solution}

\begin{problem}
已知圆 \(M: (x + \cos \theta)^2 + (y - \sin \theta)^2 = 1\)，直线 \(l: y = kx\)，下面四个命题：

\begin{circlelist}
\item 对任意实数 \(k\) 和 \(\theta\)，直线 \(l\) 和圆 \(M\) 相切；
\item 对任意实数 \(k\) 和 \(\theta\)，直线 \(l\) 和圆 \(M\) 有公共点；
\item 对任意实数 \(\theta\)，必存在实数 \(k\)，使得直线 \(l\) 和圆 \(M\) 相切；
\item 对任意实数 \(k\)，必存在实数 \(\theta\)，使得直线 \(l\) 和圆 \(M\) 相切.
\end{circlelist}

其中真命题的代号是\fillinblank{}（写出所有真命题的代号）.
\end{problem}

\begin{answer}
B，D
\end{answer}

\begin{solution}
圆心为 \(M(-\cos\theta,\sin\theta)\)，且原点在圆上. 直线 \(l:y=kx\) 也经过原点，所以它总与圆有公共点，第二个命题为真.

直线 \(l\) 与圆相切时，因公共点只能是原点，故必须是圆在原点处的切线. 圆心到原点的半径方向为 \((\cos\theta,-\sin\theta)\)，所以切线方向可取 \((\sin\theta,\cos\theta)\). 对任意 \(k\)，取 \(\sin\theta=1/\sqrt{1+k^2}\)、\(\cos\theta=k/\sqrt{1+k^2}\)，则 \((\sin\theta,\cos\theta)\) 与直线 \(y=kx\) 的方向向量 \((1,k)\) 同向，从而存在 \(\theta\) 使 \(l\) 相切，第四个命题为真.

第一个命题不真，例如 \(k=0\)，\(\theta=0\) 时，直线 \(y=0\) 是圆的直径所在直线而非切线. 第三个命题不真，因为 \(\theta=0\) 时圆在原点处的切线为竖直直线 \(x=0\)，不能写成 \(y=kx\)（其中 \(k\) 为实数）. 所以真命题的代号为 \(B\)，\(D\).
\end{solution}

\section{解答题}
\begin{problem}
已知函数 \( f(x) = x^3 + ax^2 + bx + c \) 在 \( x = -\frac{2}{3} \) 与 \( x = 1 \) 时都取得极值.

\begin{enumerate}
\item 求 \(a\)，\(b\) 的值与函数 \(f(x)\) 的单调区间；
\item 若对 \(x \in [-1, 2]\)，不等式 \(f(x) < c^2\) 恒成立，求 \(c\) 的取值范围.
\end{enumerate}
\end{problem}

\begin{solution}
\begin{enumerate}
\item 由极值点条件，\(f'(-2/3)=f'(1)=0\). 又 \(f'(x)=3x^2+2ax+b\) 是二次函数，故
\(f'(x)=3\left(x+\frac23\right)(x-1)=3x^2-x-2\). 比较系数得 \(a=-\frac12\)，\(b=-2\).

由 \(f'(x)\) 的符号可知，\(f'(x)>0\) 在 \((-\infty,-2/3)\cup(1,+\infty)\) 上成立，\(f'(x)<0\) 在 \((-2/3,1)\) 上成立. 因此 \(f\) 在 \((-\infty,-2/3)\) 和 \((1,+\infty)\) 上单调递增，在 \((-2/3,1)\) 上单调递减.
\item 写 \(f(x)=g(x)+c\)，其中 \(g(x)=x^3-\frac12x^2-2x\). 在 \([-1,2]\) 上，关键点为 \(-2/3\)，\(1\)，并且
\(g(-1)=1/2\)，\(g(-2/3)=22/27\)，\(g(1)=-3/2\)，\(g(2)=2\). 结合单调区间，得 \(g(x)\) 在该区间上的最大值为 \(2\)，在 \(x=2\) 处取得.

因此 \(f(x)<c^2\) 对所有 \(x\in[-1,2]\) 恒成立，当且仅当 \(2+c<c^2\). 解得 \((c-2)(c+1)>0\)，所以 \(c<-1\) 或 \(c>2\).
\end{enumerate}
\end{solution}

\begin{problem}
某商场举行抽奖促销互动，抽奖规则是：从装有 9 个白球，1 个红球的箱子中每次随机地摸出一个球，记下颜色后放回. 摸出一个红球可获得奖金 10 元；摸出两个红球可获得奖金 50 元. 现有甲，乙两位顾客，规定：甲摸一次，乙摸两次. 令 \(\xi\) 表示甲，乙两人摸球后获得的奖金总额. 求：

\begin{enumerate}
\item \(\xi\) 的分布列；
\item \(\xi\) 的数学期望.
\end{enumerate}
\end{problem}

\begin{solution}
\begin{enumerate}
\item 每次摸到红球的概率为 \(1/10\)，且各次摸球相互独立. 甲不摸到红球、乙不摸到红球时 \(\xi=0\)，所以 \(P(\xi=0)=(9/10)^3=729/1000\).

\(\xi=10\) 有两种互斥情形：甲摸到红球且乙两次均未摸到红球，或甲未摸到红球且乙恰摸到一次红球. 因此 \(P(\xi=10)=\frac1{10}(\frac9{10})^2+\frac9{10}\mathrm C_2^1\frac1{10}\frac9{10}=243/1000\).

\(\xi=20\) 只能是甲摸到红球且乙恰摸到一次红球，故 \(P(\xi=20)=\frac1{10}\mathrm C_2^1\frac1{10}\frac9{10}=18/1000\).

\(\xi=50\) 只能是甲未摸到红球且乙两次均摸到红球，故 \(P(\xi=50)=\frac9{10}(\frac1{10})^2=9/1000\).

\(\xi=60\) 只能是甲和乙的两次摸球均摸到红球，故 \(P(\xi=60)=(1/10)^3=1/1000\). 除上述五种取值外概率均为零，这就得到分布列.
\item 由数学期望定义，
\(E\xi=0\cdot\frac{729}{1000}+10\cdot\frac{243}{1000}+20\cdot\frac{18}{1000}+50\cdot\frac9{1000}+60\cdot\frac1{1000}=\frac{3300}{1000}=3.3\).
\end{enumerate}
\end{solution}

\begin{problem}
如图，已知 \(\triangle ABC\) 是边长为 1 的正三角形，\(M\)，\(N\) 分别是边 \(AB\)，\(AC\) 上的点，线段 \(MN\) 经过 \(\triangle ABC\) 的中心 \(G\). 设 \(\angle MGA = \alpha\)（\(\frac{\pi}{3} \leqslant \alpha \leqslant \frac{2\pi}{3}\)）.

\begin{enumerate}
\item 试将 \(\triangle AGM\) 和 \(\triangle AGN\) 的面积（分别记为 \(S_1\) 与 \(S_2\)）表示为 \(\alpha\) 的函数；
\item 求 \(y = \frac{1}{S_1^2} + \frac{1}{S_2^2}\) 的最大值与最小值.
\end{enumerate}

\bitmapfigure[max width=0.25\linewidth]{img/2006/jiangxi_science/q19_fig1.png}
\end{problem}

\begin{solution}
\begin{enumerate}
\item 正三角形的中心到顶点的距离为 \(AG=\sqrt{3}/3\)，且 \(\angle GAB=\angle GAC=\pi/6\). 在三角形 \(AGM\) 中，\(\angle MGA=\alpha\)，所以 \(\angle AMG=5\pi/6-\alpha\). 由正弦定理，
\(GM=\dfrac{AG\sin(\pi/6)}{\sin(5\pi/6-\alpha)}=\dfrac{AG}{2\sin(\alpha+\pi/6)}\). 因此
\(S_1=\dfrac12AG\cdot GM\sin\alpha=\dfrac{\sin\alpha}{12\sin(\alpha+\pi/6)}\).

直线 \(MN\) 经过 \(G\)，故 \(\angle AGN=\pi-\alpha\). 在三角形 \(AGN\) 中，\(\angle ANG=\alpha-\pi/6\)，同理得 \(GN=\dfrac{AG}{2\sin(\alpha-\pi/6)}\)，从而
\(S_2=\dfrac12AG\cdot GN\sin(\pi-\alpha)=\dfrac{\sin\alpha}{12\sin(\alpha-\pi/6)}\).
\item 由上式，
\(y=\dfrac{144}{\sin^2\alpha}\left[\sin^2(\alpha+\pi/6)+\sin^2(\alpha-\pi/6)\right]\). 利用 \(\sin^2u+\sin^2v=1-\cos(u+v)\cos(u-v)\)，得
\(y=\dfrac{144(1-\frac12\cos2\alpha)}{\sin^2\alpha}=144\dfrac{2-\cos2\alpha}{1-\cos2\alpha}\).

当 \(\pi/3\leqslant\alpha\leqslant2\pi/3\) 时，\(-1\leqslant\cos2\alpha\leqslant-1/2\). 函数 \((2-u)/(1-u)\) 在此区间上递增，所以 \(y\) 在 \(\cos2\alpha=-1\)，即 \(\alpha=\pi/2\) 时取得最小值 \(216\)；在 \(\cos2\alpha=-1/2\)，即 \(\alpha=\pi/3\) 或 \(2\pi/3\) 时取得最大值 \(240\).
\end{enumerate}
\end{solution}

\begin{problem}
如图，在三棱锥 \(A - BCD\) 中，侧面 \(ABD\)，\(ACD\) 是全等的直角三角形，\(AD\) 是公共的斜边，且 \(AD = \sqrt{3}\)，\(BD = CD = 1\)，另一个侧面 \(ABC\) 是正三角形.

\begin{enumerate}
\item 求证：\(AD \perp BC\)；
\item 求二面角 \(B - AC - D\) 的大小；
\item 在直线 \(AC\) 上是否存在一点 \(E\)，使 \(ED\) 与面 \(BCD\) 成 \(30^{\circ}\) 角？若存在，确定 \(E\) 的位置；若不存在，说明理由.
\end{enumerate}

\bitmapfigure[max width=0.42\linewidth]{img/2006/jiangxi_science/q20_fig1.png}
\end{problem}

\begin{solution}
\begin{enumerate}
\item 由两个直角三角形全等且 \(AD\) 为公共斜边，得 \(AB=AC=\sqrt{AD^2-CD^2}=\sqrt2\). 设 \(\overrightarrow{DA}=\bs{u}\)，\(\overrightarrow{DB}=\bs{v}\)，\(\overrightarrow{DC}=\bs{w}\). 由于 \(\angle ABD=\angle ACD=90^\circ\)，有
\((\bs{u}-\bs{v})\cdot(-\bs{v})=0\)，\((-\bs{w})\cdot(\bs{u}-\bs{w})=0\)，从而 \(\bs{u}\cdot\bs{v}=\bs{u}\cdot\bs{w}=1\). 因此
\(\overrightarrow{AD}\cdot\overrightarrow{BC}=(-\bs{u})\cdot(\bs{w}-\bs{v})=0\)，故 \(AD\perp BC\).
\item 建立空间直角坐标系，取
\(A=(0,0,0)\)，\(C=(\sqrt2,0,0)\)，\(B=(\sqrt2/2,\sqrt6/2,0)\)，\(D=(\sqrt2,\sqrt6/3,\sqrt3/3)\). 这些坐标满足题设的边长和直角条件.

平面 \(ABC\) 内垂直于 \(AC\) 的单位向量可取 \(\bs{u}_1=(0,1,0)\)，平面 \(ACD\) 内垂直于 \(AC\) 的单位向量可取 \(\bs{u}_2=(0,\sqrt6/3,\sqrt3/3)\). 所以二面角 \(B-AC-D\) 的余弦为 \(\bs{u}_1\cdot\bs{u}_2=\sqrt6/3\)，其大小为 \(\arccos(\sqrt6/3)\).
\item 在上述坐标系中，平面 \(BCD\) 的一个单位法向量为 \(\bs{n}=(-\sqrt2/2,-\sqrt6/6,\sqrt3/3)\). 设 \(E=(x,0,0)\) 在直线 \(AC\) 上，则
\(\overrightarrow{ED}=(\sqrt2-x,\sqrt6/3,\sqrt3/3)\)，\(|\overrightarrow{ED}|^2=(x-\sqrt2)^2+1\)，且 \(\overrightarrow{ED}\cdot\bs{n}=x/\sqrt2-1\).

设直线 \(ED\) 与平面 \(BCD\) 所成的角为 \(\varphi\)，则
\(\sin\varphi=\dfrac{|\overrightarrow{ED}\cdot\bs{n}|}{|\overrightarrow{ED}|}=\dfrac{|x-\sqrt2|/\sqrt2}{\sqrt{(x-\sqrt2)^2+1}}\). 令 \(\varphi=\pi/6\)，平方化简得 \((x-\sqrt2)^2=1\). 因此存在两个点：\(x=\sqrt2-1\) 和 \(x=\sqrt2+1\)，即它们分别位于点 \(C\) 两侧且均满足 \(CE=1\).
\end{enumerate}
\end{solution}

\begin{problem}
如图，椭圆 \(Q: \frac{x^2}{a^2} + \frac{y^2}{b^2} = 1\)（\(a > b > 0\)）的右焦点为 \(F(c, 0)\)，过点 \(F\) 的一动直线 \(m\) 绕点 \(F\) 转动，并且交椭圆于 \(A\)，\(B\) 两点，\(P\) 为线段 \(AB\) 的中点.

\begin{enumerate}
\item 求点 \(P\) 的轨迹 \(H\) 的方程；
\item 在 \(Q\) 的方程中，令 \(a^2 = 1 + \cos \theta + \sin \theta\)，\(b^2 = \sin \theta\)（\(0 < \theta \leqslant \frac{\pi}{2}\)），确定 \(\theta\) 的值，使原点距椭圆的右准线 \(l\) 最远，此时，设 \(l\) 与 \(x\) 轴交点为 \(D\)，当直线 \(m\) 绕点 \(F\) 转动到什么位置时，三角形 \(ABD\) 的面积最大？
\end{enumerate}

\bitmapfigure[max width=0.42\linewidth]{img/2006/jiangxi_science/q21_fig1.png}
\end{problem}

\begin{solution}
\begin{enumerate}
\item 设 \(A(x_1,y_1)\)、\(B(x_2,y_2)\)，\(P(x,y)\)，则 \(x=(x_1+x_2)/2\)，\(y=(y_1+y_2)/2\). 由 \(A\)，\(B\) 在椭圆上，得
\(b^2x_1^2+a^2y_1^2=a^2b^2\)，\(b^2x_2^2+a^2y_2^2=a^2b^2\). 两式相减，并利用直线 \(AB\) 经过 \(F(c,0)\)，得
\(b^2(x_1+x_2)+a^2\dfrac{y_1-y_2}{x_1-x_2}(y_1+y_2)=0\)，其中非竖直情形下 \(\dfrac{y_1-y_2}{x_1-x_2}=\dfrac{y}{x-c}\). 代入中点坐标并化简，得到
\(b^2x(x-c)+a^2y^2=0\)，即 \(b^2x^2+a^2y^2-b^2cx=0\). 当 \(AB\) 垂直于 \(x\) 轴时，\(P=(c,0)\)，同样满足该方程. 因此点 \(P\) 的轨迹为
\(H:b^2x^2+a^2y^2-b^2cx=0\).
\item 椭圆右准线方程为 \(x=a^2/c\)，所以原点到右准线的距离为 \(a^2/c\). 由 \(c^2=a^2-b^2=1+\cos\theta\)，且 \(c>0\)，有
\(\dfrac{a^2}{c}=\dfrac{1+\cos\theta+\sin\theta}{\sqrt{1+\cos\theta}}=\sqrt2\left(\cos\dfrac\theta2+\sin\dfrac\theta2\right)\).

因为 \(0<\theta/2\leqslant\pi/4\)，上式在 \(\theta/2=\pi/4\)，即 \(\theta=\pi/2\) 时取得最大值. 此时 \(a^2=2\)，\(b^2=1\)，\(c=1\)，椭圆为 \(x^2/2+y^2=1\)，右准线为 \(x=2\)，故 \(D=(2,0)\).

设非竖直直线 \(m\) 为 \(y=k(x-1)\)，与椭圆交于 \(A(x_1,y_1)\)、\(B(x_2,y_2)\). 代入椭圆方程得 \((1+2k^2)x^2-4k^2x+2k^2-2=0\)，所以
\(|x_1-x_2|=\dfrac{2\sqrt{2}\sqrt{1+k^2}}{1+2k^2}\). 又 \(|AB|=\sqrt{1+k^2}|x_1-x_2|\)，点 \(D\) 到直线 \(m\) 的距离为 \(|k|/\sqrt{1+k^2}\)，故
\(S_{\triangle ABD}=\dfrac12|AB|\,d(D,m)=\dfrac{\sqrt2|k|\sqrt{1+k^2}}{1+2k^2}\). 令 \(t=k^2\)，则
\(S_{\triangle ABD}^{2}=\dfrac{2t(t+1)}{(1+2t)^2}\)，其导数为 \(\dfrac{2}{(1+2t)^3}>0\). 因此在非竖直情形下，面积的平方随 \(t\) 增大而增大，且 \(\displaystyle\lim_{t\to\infty}S_{\triangle ABD}^{2}=\frac12\). 当 \(m\) 为竖直直线 \(x=1\) 时，椭圆交点为 \((1,\pm1/\sqrt2)\)，故 \(|AB|=\sqrt2\)，\(d(D,m)=1\)，三角形面积为 \(\sqrt2/2\)，其平方为 \(1/2\). 所以面积最大时 \(m\perp x\) 轴.
\end{enumerate}
\end{solution}

\begin{problem}
已知数列 \(\{a_{n}\}\) 满足：\(a_{1} = \frac{3}{2}\)，且 \(a_{n} = \frac{3na_{n-1}}{2a_{n-1} + n - 1}\)（\(n \geqslant 2\)，\(n \in \mathbf{N}^{*}\)）.

\begin{enumerate}
\item 求数列 \(\{a_{n}\}\) 的通项公式；
\item 证明：对一切正整数 \(n\)，不等式 \(a_1 \cdot a_2 \cdot \cdots \cdot a_n < 2 \cdot n!\) 恒成立.
\end{enumerate}
\end{problem}

\begin{solution}
\begin{enumerate}
\item 由递推式取倒数，得
\(\dfrac1{a_n}=\dfrac2{3n}+\dfrac{n-1}{3n}\cdot\dfrac1{a_{n-1}}\). 令 \(b_n=\dfrac n{a_n}\)，则 \(b_1=\dfrac23\)，且
\(b_n=\dfrac23+\dfrac13b_{n-1}\). 于是 \(b_n-1=\dfrac13(b_{n-1}-1)\)，从而 \(b_n-1=-3^{-n}\)，即 \(b_n=1-3^{-n}\). 所以
\(a_n=\dfrac{n}{1-3^{-n}}=\dfrac{n3^n}{3^n-1}\).
\item 由通项公式，
\(a_1a_2\cdots a_n=n!\prod_{k=1}^n\dfrac1{1-3^{-k}}\). 对 \(0\leqslant x_k<1\)，用归纳法证明
\[
\prod_{k=1}^{n+1}(1-x_k)
\geqslant\left(1-\sum_{k=1}^{n}x_k\right)(1-x_{n+1})
=1-\sum_{k=1}^{n+1}x_k+x_{n+1}\sum_{k=1}^{n}x_k
\geqslant1-\sum_{k=1}^{n+1}x_k
\]
其中初始情形 \(n=1\) 成立，归纳步利用 \(1-x_{n+1}\geqslant0\). 取 \(x_k=3^{-k}\)，得到
\(\prod_{k=1}^n(1-3^{-k})\geqslant1-\sum_{k=1}^n3^{-k}=\dfrac12+\dfrac1{2\cdot3^n}>\dfrac12\). 因此
\(a_1a_2\cdots a_n=n!\left[\prod_{k=1}^n(1-3^{-k})\right]^{-1}<2n!\)，命题得证.
\end{enumerate}
\end{solution}
